\section{Recomendaciones}
\label{sec:recomendaciones}

Para que el sistema siga funcionando de forma segura y estable, se recomienda:

\subsection{Recomendaciones técnicas}

\begin{itemize}
    \item \textbf{Cuidar el servidor:} Se aconseja guardar memoria de las páginas más usadas y ordenar los datos que más se buscan (como el carnet y la especialidad) para que la página siga rápida.
    \item \textbf{Guardar copias siempre:} Hacer copias automáticas todos los días en tiempo de matrícula y guardarlas en dos lugares, para no perder nada.
    \item \textbf{Mantener todo actualizado:} Mantener actualizados los programas del sistema para cuidar la seguridad y que nadie entre sin permiso ni lo tumbe.
    \item \textbf{Tablero Semestral de Indicadores:} Implementar un tablero de monitoreo semestral de los indicadores del sistema (espera menor a 1.2 s, 0 cupos repetidos, papeles completos 100 \%, facilidad de uso de al menos 90 \%, errores menores a 5 \%) a cargo del tutor guía y la secretaría académica, con reporte a la Dirección del MTP y a la Dirección Distrital para la mejora continua.
\end{itemize}

\subsection{Recomendaciones institucionales}

\begin{itemize}
    \item \textbf{Campañas Preventivas de Difusión:} La Dirección del MTP San Julián debe coordinar con las unidades educativas de origen la difusión de guías gráficas y videos cortos instructivos sobre el uso de la plataforma antes de la apertura oficial de cupos.
    \item \textbf{Garantía de Conectividad Institucional:} Gestionar ante el Gobierno Autónomo Municipal de San Julián una línea de internet dedicada de fibra óptica para el servidor institucional durante los períodos de alta demanda.
    \item \textbf{Continuidad del Software:} Fomentar que futuros estudiantes de la carrera de SISTEMAS INFORMÁTICOS asuman como proyectos de grado las fases de escalamiento del sistema (notificaciones por WhatsApp y gestión académica integral).
    \item \textbf{Capacitación de Secretaría Académica:} Programar talleres de práctica para revisar papeles y recuperar contraseñas en menos de 2 minutos con los nuevos estudiantes, para que el proceso siga en orden y sea justo para todos.
\end{itemize}
